% Deck figure n09: why a hash key moves terms and a conserved-parity key does not (README Sections 2.5 and 3.2).
% The same four GPUs and the same branching gate (generator Q) under two keys.
%   left    key = hash of the whole Pauli string (a fixed distribution function, as in the CPU implementations of
%           Section 2.5): the gate turns P into P and P + Q, the hash of P + Q differs from the hash of P, and
%           the new term is sent to the GPU that owns its key. This happens at every branching gate.
%   right   key = the parities that all gates of the window keep (r(Q) = 0 for every generator Q of the window):
%           r(P + Q) = r(P), so the new term and the term it merges with are in the sector of P.
% The blocks are a sketch: gray or sector-colored = terms, the block with the red frame = the new term.
% No measured numbers in this figure.
\begin{tikzpicture}[x=1cm, y=1cm, font=\rmfamily\footnotesize]
  \tikzset{tm/.style={draw=black!70, line width=0.35pt, inner sep=0pt, minimum width=1.04cm, minimum height=0.42cm,
                      font=\rmfamily\footnotesize},
           ts/.style={tm, minimum width=0.92cm}}
  % ================= left: hash of the whole Pauli string =================
  \node[frame, anchor=north west, minimum width=6.7cm, minimum height=7.1cm] at (0,0) {};
  \node[lttl, anchor=north] at (3.35,-0.05) {Key = Hash of the Whole Pauli String};
  \node[blk, fill=genc, anchor=north west, minimum width=6.3cm, minimum height=0.5cm, font=\rmfamily\footnotesize] at (0.2,-0.8)
       {A Branching Gate With Generator $Q$};
  \foreach \g/\x in {0/0.2, 1/1.82, 2/3.44, 3/5.06}
  {
    \draw[sarr] ({\x+0.72},-1.3) -- ({\x+0.72},-1.66);
    \node[gpu, anchor=north west, minimum width=1.44cm, minimum height=3.15cm] at (\x,-1.68) {};
    \node[cttl, anchor=north] at ({\x+0.72},-1.7) {GPU \g};
  }
  \foreach \x/\y in {0.92/-2.45, 0.92/-3.0, 0.92/-3.55, 2.54/-2.45, 2.54/-3.0, 4.16/-2.45, 4.16/-3.0,
                     5.78/-2.45, 5.78/-3.0, 5.78/-3.55}
    \node[tm, fill=black!22] at (\x,\y) {};
  \node[tm, fill=white] at (2.54,-3.55) {$P$};
  \node[tm, fill=white, draw=dred, line width=0.9pt] at (5.78,-4.1) {$P \oplus Q$};
  \draw[arr, dred] (2.54,-3.77) -- (2.54,-5.25) -- (5.78,-5.25) -- (5.78,-4.33);
  \node[note, text=dred, anchor=north] at (4.16,-5.3) {Sent at Every Branching Gate};
  \node[note] at (3.35,-6.25) {The Gate Changes the String,};
  \node[note] at (3.35,-6.58) {So the Hash Changes};
  % ================= right: the parities that the gates of the window keep =================
  \node[frame, anchor=north west, minimum width=6.7cm, minimum height=7.1cm] at (7.2,0) {};
  \node[lttl, anchor=north] at (10.55,-0.05) {Key = Parities the Window's Gates Keep};
  \node[blk, fill=genc, anchor=north west, minimum width=6.3cm, minimum height=0.5cm, font=\rmfamily\footnotesize] at (7.4,-0.8)
       {A Branching Gate With Generator $Q$};
  \foreach \g/\x/\f in {0/7.4/sa, 1/9.02/sb, 2/10.64/sc, 3/12.26/sd}
  {
    \draw[sarr] ({\x+0.72},-1.3) -- ({\x+0.72},-1.66);
    \node[gpu, anchor=north west, minimum width=1.44cm, minimum height=3.15cm] at (\x,-1.68) {};
    \node[cttl, anchor=north] at ({\x+0.72},-1.7) {GPU \g};
    \node[draw=black!70, line width=0.45pt, fill=\f!45, anchor=north west, minimum width=1.2cm, minimum height=2.57cm] at ({\x+0.12},-2.14) {};
    \node[ittl, anchor=north] at ({\x+0.72},-2.16) {Sector \g};
  }
  \foreach \x/\y/\f in {8.12/-2.95/sa, 8.12/-3.55/sa, 8.12/-4.15/sa, 11.36/-2.95/sc, 11.36/-3.55/sc,
                        12.98/-2.95/sd, 12.98/-3.55/sd, 12.98/-4.15/sd}
    \node[ts, fill=\f] at (\x,\y) {};
  \node[ts, fill=sb] at (9.74,-2.95) {$P$};
  \node[ts, fill=sb, draw=dred, line width=0.9pt] at (9.74,-4.15) {$P \oplus Q$};
  \draw[sarr] (9.74,-3.17) -- (9.74,-3.92);
  \node[key] at (10.55,-5.47) {No Term Leaves Its Sector During the Window};
  \node[note] at (10.55,-6.25) {The New Term and Its Merge Partner};
  \node[note] at (10.55,-6.58) {Stay in the Same Sector};
\end{tikzpicture}
